\documentclass[11pt]{article}
\usepackage[a4paper,margin=27mm]{geometry}
\usepackage{amsmath,amssymb,amsthm,mathtools,bm}
\usepackage[T1]{fontenc}
\usepackage{lmodern,microtype,xcolor}
\usepackage[colorlinks=true,linkcolor=blue!45!black,citecolor=blue!45!black,urlcolor=blue!45!black]{hyperref}
\numberwithin{equation}{section}
\newtheorem{theorem}{Theorem}[section]
\newtheorem{lemma}[theorem]{Lemma}
\newtheorem{proposition}[theorem]{Proposition}
\newtheorem{corollary}[theorem]{Corollary}
\theoremstyle{remark}\newtheorem{remark}[theorem]{Remark}
\DeclareMathOperator{\Pf}{Pf}
\DeclareMathOperator{\dist}{dist}
\DeclareMathOperator{\supp}{supp}
\newcommand{\ii}{\mathrm i}
\newcommand{\dd}{\mathrm d}
\newcommand{\Rea}{\operatorname{Re}}
\newcommand{\Ima}{\operatorname{Im}}
\newcommand{\T}{\mathbb T}
\newcommand{\X}{\mathcal X}
\newcommand{\Mtwo}{\mathcal M^2}
\allowdisplaybreaks[2]
\title{The full low-temperature Ising susceptibility:\\ a regular-coordinate proof of the natural boundary}
\author{}
\date{}
\begin{document}
\maketitle

\begin{abstract}
We prove that the complete connected bulk susceptibility of the
isotropic square-lattice Ising model, continued from its low-temperature
pure phase, has the unit circle as a natural boundary in
$s=\sinh(2\beta_{\rm phys}J)$. An exact lattice Fourier sum expresses
each full coefficient as an angular integral with a total-angle
constraint and one product denominator. At a dense family of prime
roots of unity, the first singular even-particle term has a nonzero
$\varepsilon^{-1/2}$ derivative, computed directly on this constraint.
An elementary four-point determinant identity, propagated by
Pfaffian condensation, determines the exact local coefficient. A deformation preserving the total angle and one regular
coordinate on its one-sided part give
an absolute logarithmic bound for the entire differentiated
higher-particle tail. One collective tangential field treats all right-side configurations.
Oblique coarea and one mean scale control every collision diameter,
including diameters below the radial distance; a bounded anchor
field handles the remaining regions with a separated product kernel. The same geometry
controls the finite-order complement of the first singular chart. These estimates establish
noncancellation for the full infinite expansion.
\end{abstract}

\section{The bulk function and the argument}

Consider the zero-field Ising model on the infinite square lattice,
with equal positive nearest-neighbor coupling $J$. Set
\[
 s=\sinh(2\beta_{\rm phys}J),\qquad
 \Mtwo(s)=(1-s^{-4})^{1/4},
\]
where the branch of $\Mtwo$ tends to one at infinity.
For real $s>1$ the connected response in the plus pure phase is
\begin{equation}\label{eq:physical}
 \X(s)=\beta_{\rm phys}^{-1}\chi(s)
 =\sum_{(a,b)\in\mathbb Z^2}
       \bigl(\langle\sigma_{00}\sigma_{ab}\rangle_+
                     -\Mtwo(s)\bigr).
\end{equation}
Spin reversal gives the same function in the minus pure phase.
The physical inputs are Yang's magnetization formula and the
low-temperature correlation expansion of Palmer--Tracy
\cite{Yang,PalmerTracy}, in the normalization recorded by
Tracy--Widom \cite[Appendix 1, pp.~1132--1133]{TracyWidom}.

\begin{theorem}\label{thm:main}
The exterior holomorphic function $\X$ determined by
\eqref{eq:physical} has the entire circle $|s|=1$ as a natural
boundary. The same local noncontinuation holds for the physical
susceptibility $\chi=\beta_{\rm phys}\X$, with
$\beta_{\rm phys}(s)=\operatorname{arsinh}(s)/(2J)$ continued from
the real low-temperature interval.
\end{theorem}

We use coefficients $C_N$ that already contain the full lattice sum:
\[
 \X=\Mtwo\sum_{\substack{N\ge2\\N\ \mathrm{even}}} C_N.
\]
At a prime root $\zeta=e^{2\pi ia/p}$ in the first quadrant,
the first singular even order is $N_*=2p$. For
$k=N_*^2/2-1$ and $s_\varepsilon=(1+\varepsilon)\zeta$, we prove
\begin{equation}\label{eq:two-goals}
 \begin{split}
 C_{N_*}^{(k)}(s_\varepsilon)
    &=\Lambda_\alpha\varepsilon^{-1/2}+o(\varepsilon^{-1/2}),
                  \qquad \Lambda_\alpha\ne0,\\
 \sum_{\substack{N>N_*\\N\ \mathrm{even}}}
          |C_N^{(k)}(s_\varepsilon)|
    &=O(1+|\log\varepsilon|).
 \end{split}
\end{equation}
All constants may depend on this fixed point and derivative order.
The first estimate is a local shape integral. For the second, we
deform the complete constrained angular torus and then partition it.
The regular coordinate includes the moving branch center, so
parameter differentiation has uniform coefficients even at a
dispersion branch. One collective field freezes the angular
constraint and the total phase on the entire right part. The
remaining left-root part has a separated product kernel and needs
only a bounded anchor field. Determinant minors retain factorial
decay, while oblique coarea gives one estimate for every right-side
diameter. A fixed finite mean gap places the small-order collision
charts away from the branch.

\section{The full angular representation}
\label{sec:representation}\label{sec:angular-bulk}

Write
\[
 S(s)=s+s^{-1},\qquad
 W(y;s)=S(s)-\tfrac12(y+y^{-1}),\qquad
 z=W-\sqrt{W^2-1}=e^{-\gamma}.
\]
For $|s|>1$ and $|y|=1$, choose the unique interior root $|z|<1$,
continued from real $s>1$. Indeed $S(s)\notin[-2,2]$ in the
exterior, so $W$ cannot lie in $[-1,1]$. The local logarithms
$\gamma$ have positive real part, and
$\sinh\gamma=(z^{-1}-z)/2$ is single-valued.

On $y_i=e^{i\theta_i}$ put $Z=\prod_i z_i$ and
\begin{equation}\label{eq:angular-kernel}
 K_{il}=\frac{\sin((\theta_i-\theta_l)/2)}
                  {\sinh((\gamma_i+\gamma_l)/2)},
 \qquad K_{ii}=0.
\end{equation}
Changing a half-angle or logarithmic sign changes one row and its
corresponding column, so the determinant is unambiguous and periodic.
Each kernel entry has the rational square
\begin{equation}\label{eq:kernel-square}
 K_{il}^2=-\frac{(y_i-y_l)^2z_i z_l}
                    {y_i y_l(1-z_i z_l)^2}.
\end{equation}
The determinant is used directly in the physical expansion.
Section~\ref{sec:determinant-local} proves its collision zeros
and the precise local coefficient needed below.

Let the constrained torus and its normalized Haar measure be
\[
 \mathbb H_N=\{\theta\in\mathbb T^N:
                      \textstyle\sum_i\theta_i=0\pmod{2\pi}\},
\]
\[
 d\mu_N=\prod_{i=1}^{N-1}\frac{d\theta_i}{2\pi},
 \qquad \theta_N=-\sum_{i<N}\theta_i\pmod{2\pi}.
\]
Equivalently, $d\mu_N$ is Euclidean hypersurface measure divided
by $\sqrt N(2\pi)^{N-1}$.

\begin{proposition}[Complete lattice sum]\label{prop:angular-bulk}
For even $N\ge2$, define
\begin{equation}\label{eq:angular-bulk}
 C_N(s)=\frac1{N!}\int_{\mathbb H_N}
       \frac{1+Z}{1-Z}\det K\prod_{i=1}^N\frac1{\sinh\gamma_i}
                         \,d\mu_N.
\end{equation}
Then the exterior continuation of the physical susceptibility is
\begin{equation}\label{eq:bulk}
 \X(s)=\Mtwo(s)\sum_{\substack{N\ge2\\N\ \mathrm{even}}} C_N(s).
\end{equation}
Every fixed differentiated series converges normally on compact
exterior parameter sets. All lattice sites, including the origin,
are included in \eqref{eq:bulk}.
\end{proposition}
\begin{proof}
For real $s>1$, the physical input is the Fredholm formula
\[
 \langle\sigma_{00}\sigma_{ab}\rangle_+
       =\Mtwo\det(I+g_{a,|b|}),\qquad
 g_{a,|b|}(\theta_i,\theta_l)
       =e^{ia\theta_i-|b|\gamma_i}K_{il},
\]
with measure $d\theta/(2\pi\sinh\gamma)$.
The originally stated kernel is
$\sinh((\gamma_i-\gamma_l)/2)/\sin((\theta_i+\theta_l)/2)$.
Subtracting the dispersion equations and using the difference
formulas for cosine and hyperbolic cosine gives exactly
\eqref{eq:angular-kernel}, with removable values filled in by
continuity. Each Fredholm minor is the product of its row weights
and $\det K$. The odd minors vanish because $K$ is skew.
The expansion therefore gives
\[
 \langle\sigma_{00}\sigma_{ab}\rangle_+
   =\Mtwo\left(1+\sum_{\substack{N\ge2\\N\ \mathrm{even}}}
                         c_N(a,b;s)\right),
\]
where
\[
 c_N(a,b;s)=\frac1{N!}\int_{\mathbb T^N}
       e^{ia\sum_i\theta_i}Z^{|b|}\det K
                    \prod_i\frac{d\theta_i}{2\pi\sinh\gamma_i}.
\]
For this fixed $s$, all entries and one-body factors are bounded on
a common small angular strip, on which $|z_i|\le e^{-c}<1$.
Shift every angle by $i\eta\,\operatorname{sgn}(a)$, with fixed
small $\eta>0$. Periodicity and analyticity give the sitewise bound
\[
 |c_N(a,b;s)|\le\frac{C^N N^{N/2}}{N!}
                    e^{-N(\eta|a|+c|b|)}.
\]
Hadamard's inequality supplies the determinant factor. The bound is
summable in all sites and all even $N\ge2$, so the lattice and
particle sums may be interchanged.

Summing over $b\in\mathbb Z$ gives $(1+Z)/(1-Z)$. To justify
the other direction with its exact normalization, for a periodic
function $F$ on $\mathbb T^N$ define
\[
 Q(v)=\int_{\mathbb T^{N-1}}
 F(\theta_1,\ldots,\theta_{N-1},v-\textstyle\sum_{i<N}\theta_i)
                 \prod_{i<N}\frac{d\theta_i}{2\pi}.
\]
Its Fourier coefficient of index $-a$ is
$\int e^{ia\sum\theta_i}F\prod_i d\theta_i/(2\pi)$.
For the present analytic integrand the Fourier series is absolutely
convergent. Summing those coefficients gives $Q(0)$, precisely
\eqref{eq:angular-bulk}. Subtracting $\Mtwo$ in the correlation
formula removes its zeroth-order term and proves \eqref{eq:bulk}.
At the origin the full correlation is one, so its connected
contribution in this sum is exactly $1-\Mtwo$.

On a compact exterior neighborhood the unit-circle roots satisfy
$m\le|z_i|\le q<1$. The complete kernel entries, one-body factors,
and product denominator in \eqref{eq:angular-bulk} are uniformly
bounded. Since $\mu_N(\mathbb H_N)=1$, Hadamard and
$N!\ge(N/e)^N$ give
\[
 |C_N(s)|\le C_0(C_1/\sqrt N)^N.
\]
Cauchy's inequality on a slightly larger compact neighborhood gives
the same summability for every fixed derivative. Analytic continuation
from the real interval proves the exterior identity.
\end{proof}

Changing $s,y_i,z_i$ to $-s,-y_i,-z_i$ preserves the constrained
torus when $N$ is even. It changes each complete matrix entry and
each one-body factor by a minus sign; the full coefficient is
unchanged. Thus $C_N$ and $\X$ are even. Their values also commute
with complex conjugation, first on the physical germ and then by
analytic continuation.

\section{Prime roots with a unique first cosine pair}
\label{sec:prime-roots}

Choosing the two dispersion angles equal gives a dense family of roots
of unity. At the first even singular order there is a unique cosine
pair. Its two signed angular configurations are exchanged by reflection.

\begin{lemma}\label{lem:prime-root-family}
Let $p\geq5$ be prime and $1\leq a<p/4$. Put
\[
 \alpha=\frac{2\pi a}{p},\qquad
 \zeta=e^{i\alpha},\qquad S_*=2\cos\alpha,
 \qquad N_*=2p.
\]
Then the first positive even integer $n$ for which
\[
 S_*=\cos(2\pi l/n)+\cos(2\pi m/n)
\]
is $n=N_*$. At that order the unordered cosine pair is uniquely
$\{\cos\alpha,\cos\alpha\}$. These points $\zeta$ are dense in the open
upper-right quadrant of the unit circle. If
\[
 b=\arccos(S_*-1),\qquad \xi=e^{-ib},
\]
then $\xi$ is algebraic and is not a root of unity.
\end{lemma}

\begin{proof}
Write $\zeta_p=e^{2\pi i/p}$, $d=(p-1)/2$, and
$u_j=\zeta_p^j+\zeta_p^{-j}$ for $1\leq j\leq d$. In
$E=\mathbb Q(\zeta_p+\zeta_p^{-1})$, every $u_j$ has trace $-1$:
its conjugates sum the nontrivial $p$-th roots of unity once each.
In particular $S_*=u_a>0$ is irrational, since rationality would give
$d u_a=-1$.

If $n<2p$ is even, then $p\nmid n$. The compositum of
$\mathbb Q(\zeta_p)$ and $\mathbb Q(\zeta_n)$ is
$\mathbb Q(\zeta_{pn})$. Its cyclotomic degree is
$\varphi(pn)=\varphi(p)\varphi(n)$, so the two fields have intersection
$\mathbb Q$. Thus a Nickel equation of order $n$ would make $S_*$
rational. The displayed pair $(\alpha,\alpha)$ is available at order
$2p$, proving the first-order assertion.

Next $\xi+\xi^{-1}=2u_a-2$, so $\xi$ is algebraic. If $\xi$ were a
root of unity, every conjugate of $2u_a-2$ would lie in $[-2,2]$.
Its trace in $E$, however, is $-2-2d<-2d$. This proves nonresonance.

At order $2p$ a cosine equal to $-1$ is impossible, since its partner
would have to be $S_*+1>1$. A cosine equal to $1$ would make $\xi$
a $2p$-th root of unity. All remaining twice-cosines belong to
$\{\pm u_1,\ldots,\pm u_d\}$. The numbers $u_1,\ldots,u_d$ are
linearly independent over $\mathbb Q$: a relation yields a polynomial
of degree at most $p-1$, with zero constant coefficient, vanishing at
$\zeta_p$. It is a multiple of $\Phi_p(X)=1+X+\cdots+X^{p-1}$, and
its constant coefficient forces that multiple to be zero.
Consequently $2u_a=v+w$, with $v,w\in\{\pm u_j\}$, implies
$v=w=u_a$.

Finally, given $\alpha_0\in(0,\pi/2)$, choose unbounded primes $p$
and integers $a$ nearest to $p\alpha_0/(2\pi)$. Eventually
$1\leq a<p/4$, and $2\pi a/p\to\alpha_0$.
\end{proof}

\section{An elementary product identity for the correlation determinant}
\label{sec:elementary-pfaffian}\label{sec:determinant-local}

The dispersion relation makes a four-point identity an ordinary
Vandermonde determinant.  Pfaffian condensation then proves the
identity at every even order.  We give both steps, since they also
fix all local signs and collision multiplicities.

\begin{lemma}[Four points on the dispersion curve]
Let nonzero $y_i,z_i$ satisfy
$y_i+y_i^{-1}+z_i+z_i^{-1}=2S$ for $1\le i\le4$, and set
$B_{ij}=(y_i-y_j)/(1-z_i z_j)$.  Then
\[
 B_{12}B_{34}-B_{13}B_{24}+B_{14}B_{23}
 =\left(\prod_{i=1}^4\frac{z_i}{y_i}\right)
   \prod_{i<j}B_{ij}.
\]
The equality holds wherever the expressions are defined and,
after clearing denominators, at every point of the curve.
\end{lemma}
\begin{proof}
Write $a_{ij}=y_i-y_j$, $t_i=z_i+z_i^{-1}$, and
$Z_4=z_1z_2z_3z_4$.  For a partition $ij\mid kl$ of the four
indices, direct multiplication gives
\[
 \frac{\prod_{v\in\{i,j\},\,w\in\{k,l\}}(1-z_vz_w)}{Z_4}
 =R+t_it_j+t_kt_l,
 \qquad
 R=Z_4+Z_4^{-1}+2-\sum_{v<w}
          \bigl(z_vz_w+(z_vz_w)^{-1}\bigr).
\]
Multiplying the proposed identity by
$\prod_{i<j}(1-z_i z_j)$, the common $R$ term cancels because
$a_{12}a_{34}-a_{13}a_{24}+a_{14}a_{23}=0$.
The remaining left side, divided by $Z_4$, is
\[
 \begin{split}
 &a_{12}a_{34}(t_1t_2+t_3t_4)
 -a_{13}a_{24}(t_1t_3+t_2t_4)\\
 &\hspace{25mm}+a_{14}a_{23}(t_1t_4+t_2t_3)
 =\det[\,1,\ y_i,\ t_i,\ y_it_i\,]_{i=1}^4.
 \end{split}
\]
Here the last equality is a two-column Laplace expansion.
Since $t_i=2S-y_i-y_i^{-1}$, column operations give
\[
 \det[1,y_i,t_i,y_it_i]
 =\det[1,y_i,y_i^{-1},y_i^2]
 =\frac{\det[y_i,y_i^2,1,y_i^3]}{\prod_i y_i}
 =\frac{\prod_{i<j}(y_i-y_j)}{\prod_i y_i}.
\]
This is the required cleared numerator.
\end{proof}

\begin{lemma}[Pfaffian propagation]\label{lem:pf-propagation}
Suppose a skew matrix $K$ has the property that every four-index
principal submatrix satisfies
\[
 K_{ab}K_{cd}-K_{ac}K_{bd}+K_{ad}K_{bc}
   =\prod_{v<w\in\{a,b,c,d\}}K_{vw}
 \quad(a<b<c<d).
\]
Then, at every even order for which the intervening products
can be canceled, $\operatorname{Pf}K=\prod_{i<j}K_{ij}$.
For a rational or analytic family with a dense set of such
points, the identity extends to the whole family.
\end{lemma}
\begin{proof}
Let $I$ be the first $N-4$ indices and let $a,b,c,d$ be the last
four.  The elementary Pfaffian condensation identity is
\[
 \operatorname{Pf}K_{Iabcd}\operatorname{Pf}K_I
 =\operatorname{Pf}K_{Iab}\operatorname{Pf}K_{Icd}
  -\operatorname{Pf}K_{Iac}\operatorname{Pf}K_{Ibd}
  +\operatorname{Pf}K_{Iad}\operatorname{Pf}K_{Ibc}.
\]
For completeness, when $K_I=A$ is invertible, the simultaneous
row and column elimination in
$\left[\begin{smallmatrix}A&B\\-B^T&C\end{smallmatrix}\right]$
has determinant one and gives the diagonal blocks
$A$ and $C+B^TA^{-1}B$.  Pfaffian multilinearity therefore
reduces the displayed formula to the defining four-by-four
Pfaffian formula.  The identity is polynomial, so it holds
also when $A$ is singular.

Assume the product formula at all smaller even orders.
Writing $P_I=\prod_{i<j\in I}K_{ij}$, the right side factors as
\[
 P_I^2\prod_{i\in I,\,v\in\{a,b,c,d\}}K_{iv}
 \bigl(K_{ab}K_{cd}-K_{ac}K_{bd}+K_{ad}K_{bc}\bigr)
 =P_I\prod_{i<j}K_{ij}.
\]
The left side is $P_I\operatorname{Pf}K$.  Cancel $P_I$.
The orders two and four start the induction.
\end{proof}

\begin{proposition}[Complete product and local coefficient]
\label{prop:complete-product}\label{prop:local-determinant-coefficient}
Choose local square roots $r_i^2=z_i/y_i$ and set
$K_{ij}=-ir_i r_j B_{ij}$.  For every even $N$ on the
dispersion curve,
\[
 \operatorname{Pf}K=\prod_{i<j}K_{ij},\qquad
 \det K=\prod_{i<j}K_{ij}^2.
\]
The physical kernel has precisely this local form.  The
determinant identity is independent of all square-root choices.
\end{proposition}
\begin{proof}
At order four,
\[
 \Pf K=-\Bigl(\prod_i r_i\Bigr)\Pf B
       =-\Bigl(\prod_i r_i^3\Bigr)\prod_{i<j}B_{ij}
       =\prod_{i<j}K_{ij}.
\]
Thus the preceding identity is the four-point hypothesis for $K$.  Apply Lemma~\ref{lem:pf-propagation} at
generic points, where all pair entries are nonzero, then clear
denominators and continue.  This argument also covers
degenerate parameter values by continuation.  The identity
$(\operatorname{Pf}K)^2=\det K$ gives the second equality.
Changing a square root changes one row and its column by a
minus sign and leaves both determinant expressions invariant.
\end{proof}

\begin{corollary}[Fixed-order collision divisibility]
\label{lem:alternating-division}
On every regular fixed-order collision chart, each pair entry
has a factor $\theta_i-\theta_j$; permutation covariance of
the coupled contour gives the same factor when occupancies
are present.  Thus $\det K$ has the full squared Vandermonde
factor, with bounded fixed-order jets on compact regular
charts.  Near $\theta_i=-\alpha+u_i$, $s=e^{i\alpha}$, the
physical dispersion is $\gamma_i=i\alpha+O(u_i)$, and hence
\[
 K_{ij}=\frac{u_i-u_j}{2i\sin\alpha}
          +O\bigl((u_i-u_j)|u|\bigr),\qquad
 \det K=\Delta(u)^2\left[
 (2i\sin\alpha)^{-N(N-1)}+O(|u|+|s-e^{i\alpha}|)\right].
\]
The analytic quotient and every parameter derivative retain
the same Vandermonde factor.  With all one-body measures,
$1+Z\to2$, and the factor $1/N!$ included, the normalized
local coefficient is exactly
\[
 A_\alpha=\frac{2}{N!(2\pi)^{N-1}}
           2^{-N(N-1)}\sin(\alpha)^{-N^2}>0.
\]
Indeed the combined phase is
$(-i)^N(-1)^{N(N-1)/2}=i^{-N^2}=1$ for even $N$.
\end{corollary}

\section{A product-preserving deformation of the whole-bulk integral}
\label{sec:constrained-deformation}

We use the normalized coefficients $C_N$ of
\eqref{eq:angular-bulk}.  Fix a selected point $s_*$, put
$S_*=s_*+s_*^{-1}\in(0,2)$, and write
\[
 s_\varepsilon=(1+\varepsilon)s_*,\qquad
 b=\arccos(S_*-1),\qquad \xi=e^{-ib},\qquad
 H=\log(1/\varepsilon).
\]
The constants in this section may depend on the selected point and a
fixed derivative order $k$, but never on the even particle number $N$
or on $\varepsilon$.  The infinite-tail estimates concern $N\ge N_*+2$, where
$k=N_*^2/2-1$. Set $N_c=2k+4$.
We prove logarithmic absolute bounds through $k$, which are
smaller than the first singular term. The deformation and the regular transport formulas
are also used at fixed even orders $2\le N\le N_*$ in
Section~\ref{sec:first-single}.

For the regular dispersion coordinate we write
\[
 \varphi=\arccos W,\qquad z=e^{-i\varphi},\qquad
 \gamma=i\varphi.
\]
Its physical continuation lies in the fourth quadrant near either
inverse-dispersion branch.  The angular integration coordinates remain
$\theta_i$, with $\sum_i\theta_i=0\pmod {2\pi}$.

Choose a smooth periodic function $p$ supported in $(0,\pi)$,
with $0\le p\le1$ and $p=1$ on $[\delta,\pi-\delta]$.
Here $\delta>0$ is smaller than $\alpha/2$, $b/2$, and
$(2\alpha-b)/2$, and is small enough that $b$ is inside this
plateau and $1+g'(\theta)>0$ on the upper endpoint sliver.
Here $g(\theta)=\arccos(S_*-\cos\theta)$; the choices are
possible because $b<2\alpha$ and $g'(0)=0$.  Put $m(\theta)=p(-\theta)$.  Thus $p,m$ have
disjoint supports, and both inverse branches have fixed plateau
neighborhoods.  Define
\begin{equation}\label{eq:constrained-homotopy}
 \begin{aligned}
 P&=\sum_i p(\theta_i),\qquad Q=\sum_i m(\theta_i),\\
 h_i&=\frac\tau N\{P m(\theta_i)-Q p(\theta_i)\},\qquad
 y_i(\lambda)=e^{i\theta_i+\lambda h_i},\quad 0\le\lambda\le1.
 \end{aligned}
\end{equation}
The positive number $\tau$ is fixed and sufficiently small.  In
particular $\sum_i h_i=0$, so $\prod_i y_i=1$ throughout the
deformation.

\begin{lemma}[Exact constrained deformation]
\label{lem:constrained-homotopy}
The full integral defining $C_N$ is unchanged by
\eqref{eq:constrained-homotopy}.  Its angular Jacobian and its fixed
jets have bounds $C_j^N N^{C_j}$, uniformly in $N$.
\end{lemma}
\begin{proof}
For $y=e^{v+i\theta}$,
\[
 \Im W=\Im S-\sinh v\sin\theta.
\]
On the upper semicircle the motion in
\eqref{eq:constrained-homotopy} is inward, and on the lower
semicircle it is outward.  Both motions increase $\Im W$.
For the admissible exterior parameters it remains positive, so every
$z_i$ stays on the interior sheet.  Thus $1-Z$ and all the pair
denominators $1-z_i z_l$ stay nonzero.  The kernel square \eqref{eq:kernel-square} has no
$1-y_i y_l$ denominator.  The closed constrained torus therefore
undergoes a legal holomorphic homotopy.

Write $w_i=\theta_i-i h_i$, so $\sum_iw_i=\sum_i\theta_i$.
The full matrix $I-iDh$ preserves the tangent hyperplane and induces
the identity on its one-dimensional quotient.  Its full determinant
is consequently its constrained tangent determinant.  Directly,
\[
 Dh=\frac\tau N\operatorname{diag}(P m'_i-Qp'_i)
       +\frac\tau N(m\otimes p'-p\otimes m').
\]
The diagonal and the two rank-one terms have operator norms at most
$C\tau$, independently of $N$.  Hence the Jacobian has modulus at
most $C^N$, and its inverse matrix is uniformly bounded after reducing
$\tau$.  Differentiating the determinant, or its logarithmic derivative
formula, shows that fixed jets add only a polynomial in $N$.
\end{proof}

Let $\eta$ be smooth, $0\le\eta\le1$, with $\eta(t)=0$ for
$t\le1/4$ and $\eta(t)=1$ for $t\ge1/2$.  Partition the final
deformed integral with the three weights
\[
 w_F=\eta(P)\eta(Q),\qquad
 w_L=\tfrac12(1-\eta(P))(1+\eta(Q)),\qquad
 w_U=\tfrac12(1-\eta(Q))(1+\eta(P)).
\]
Their sum is one.  Reflection $\theta\mapsto-\theta$ exchanges
$P,Q$, takes $y_i$ to $y_i^{-1}$, preserves $z_i$, and changes
$K$ to $-K$. For even $N$ its determinant and the normalized
constrained integral are unchanged. Therefore
\begin{equation}\label{eq:constrained-split}
 C_N=F_N+2L_N,
\end{equation}
where $F_N$ and $L_N$ are the final integrals with weights $w_F$
and $w_L$.  All weights are imposed after the full closed homotopy.
The support of $L_N$ has $P\le1/2$; it excludes the upper inverse
branch and consists of the lower semicircle together with two small
upper endpoint slivers.

\subsection{The sector containing both branch populations}

\begin{lemma}\label{lem:constrained-selected}
For every fixed $j\ge0$ and every even $N\ge4$,
\[
 |F_N^{(j)}(s_\varepsilon)|
    \le C_j N^{3j+1}\left(\frac{C_j}{\sqrt N}\right)^N.
\]
In particular these derivatives have a uniformly bounded absolute sum.
\end{lemma}
\begin{proof}
On this support $P,Q\ge1/4$.  Both bump supports stay a fixed
distance from the endpoints $0,\pi$.  The imaginary-part formula,
and the bounded range of $W$, imply
\[
 -\log|Z|\ge c\left(N\varepsilon+\frac{\tau PQ}{N}\right)
             \ge \frac cN.
\]
Every true lower branch has outward displacement $\tau P/N$,
and every true upper branch has inward displacement $\tau Q/N$.
Their normal forms give a branch distance at least $c/N$ in
the $W$ plane.  All other coordinates stay in fixed regular
neighborhoods.  Consequently $|\partial_s\log z_i|\le C\sqrt N$
on a disk $|s-s_\varepsilon|\le c/N^3$.  Reduce $c$ once: the
variation of $\log Z$ is at most $CN^{3/2}c/N^3$, smaller than
half its preceding gap.  The true branches remain on their interior
sheets throughout this disk.

We estimate the correlation determinant directly.
Entries with both points near the same inverse branch are bounded:
the common factor $\varphi_i+\varphi_l$ cancels, leaving
$(\varphi_i-\varphi_l)$ times a bounded analytic coefficient.
Away from an opposite-branch pair, compactness gives the same bound.
Let $B$ be the number of coordinates in the two fixed true-branch
neighborhoods.  If a pair is near opposite branches, its numerator
is bounded and
\[
 |1-z_i z_l|\ge c\sqrt{(P+Q)/N},\qquad
 |K_{il}|\le C\sqrt{N/(P+Q)}\le C\sqrt{N/B}.
\]
Indeed both regular branch coordinates remain in a common sector
of width strictly less than $\pi$ (a small fixed enlargement of the
fourth quadrant suffices), while their imaginary dispersion margins are
at least $c\tau P/N$ and $c\tau Q/N$.  Each such coordinate
contributes one to $P+Q$.  There are at most $B$ exceptional
entries in any row; all remaining entries are bounded.  Every row
therefore has Euclidean norm at most $C\sqrt N$.  Hadamard gives
\[
 |\det K|\le C^N N^{N/2}
\]
uniformly on the disk.

The one-body factors have a common independent angular majorant
\[
 f(\theta)=C\left(1+
  \operatorname{dist}(\theta,\{b,-b\})^{-1/2}\right).
\]
To verify this, use the imaginary part of the branch normal form
when $|u|\le C(\varepsilon+|h|)$ and its linear real part outside
that interval.  The small parameter disk is absorbed in the positive
branch displacement.  Thus $f\in L^{4/3}(\mathbb T)\cap L^1(\mathbb T)$.
With normalized circle measure, the constrained product integral is
$f^{*N}(0)$.  Indeed, $g=f*f$ belongs to $L^2$ by Young's
inequality, and $g*g$ is bounded and continuous by
Cauchy--Schwarz and continuity of translations in $L^2$.
Further convolution with $f\in L^1$ preserves these properties.
Thus evaluation at zero is legitimate, and Young's inequalities give
\[
 f^{*N}(0)\le \|f\|_{4/3}^{4}\|f\|_1^{N-4}\le C^N,
 \qquad N\ge4.
\]
The Jacobian costs $C^N$ by Lemma~\ref{lem:constrained-homotopy}.
The remaining product kernel costs $CN$.  Division by
$N!\ge(N/e)^N$ bounds the integral on the disk by
$N(C/\sqrt N)^N$.  Cauchy's inequality proves the assertion.
\end{proof}

\section{A regular coordinate for the one-sided contour}\label{sec:global-coordinate}


The support of $L_N$ can be lifted to one angular interval
$(-\pi-\delta,\delta)$: the remaining upper plateau has zero
weight because $p(\theta_i)=1$ would imply $P>1/2$.
The possible half-angle sign changes are row and column gauge
changes, so this lift preserves the determinant and its integral.
We use $w_i=\theta_i-ih_i$ and the physical phase
$\varphi_i=\arccos(S-\cos w_i)$ throughout.

Let $N_c=2k+4$. Choose a fixed branch width, decreased further
below if needed, so that
\begin{equation}\label{eq:finite-mean-exclusion}
 \delta_B<\frac14\min_{\substack{2\le n\le N_c\\n\ \mathrm{even}}}
       \operatorname{dist}(-b,2\pi\mathbb Z/n).
\end{equation}
The minimum is positive by Lemma~\ref{lem:prime-root-family}.
The right interval will be $[-b-\delta_B,\delta]$; its
complement on the one-sided support has a compact left root.
Only this one smooth division will be needed. The displacement
strength is chosen after these fixed margins and cutoffs.

Let
\[
 w_B(s)=-\arccos(S(s)-1),\qquad
 q=\varphi+w-w_B(s).
\]
The branch of $w_B$ is fixed by $w_B(s_*)=-b$.

\begin{lemma}[Regular inverse and right-side convexity]\label{lem:inverse}
The one-sided limiting dispersion curve has a fixed complex
neighborhood on which $q$ is a regular coordinate and $w=U(s,q)$ is
analytic, with uniformly bounded fixed jets. The inverse can be chosen
on a fixed convex neighborhood $\mathcal D_R$ of the right $q$-interval,
enlarged to contain a sufficiently short negative branch segment, so
that
\[
 U(s,0)=w_B(s),\qquad U_q(s,0)=0,\qquad
 \Rea U_{qq}(s,q)\ge c>0 \quad(q\in\mathcal D_R).
\]
The right angular interval extends from $-b-\delta_B$ to a sufficiently
small positive endpoint sliver. The smallness of $\delta_B$ and the
deformation strength guarantees that its actual $q$-images lie in
$\mathcal D_R$.
\end{lemma}
\begin{proof}
On the limiting compact right arc let
$g(\theta)=\arccos(S_*-\cos\theta)$. Choose the upper endpoint small
enough that $1+g'(\theta)>c>0$. Since
\[
 g''(\theta)=-\frac{S_*(1-\cos\theta\cos g(\theta))}{\sin^3g(\theta)}<0,
\]
the coordinate $q=g(\theta)+\theta+b$ is strictly increasing. Near the
branch, use $\varphi$ itself: the lower inverse angular function is
even analytic and $w-w_B=O(\varphi^2)$, so
$q=\varphi+O(\varphi^2)$ is a regular analytic coordinate.

There is also an explicit right inverse. Put $v=q+w_B(s)$; then
$v=w+\varphi$ and
\begin{equation}\label{eq:q-1}
U(s,q)=\frac v2-
 \arccos\left(\frac{S(s)}{2\cos(v/2)}\right).
\end{equation}
The branch is fixed by $w<\varphi$ on the limiting right arc. Its real
$v$-interval is contained strictly in $(-2\alpha,2\alpha)$, after the
upper sliver is shortened. With $x=v/2$ and
$D=4\cos^2x-S_*^2>0$, differentiation gives
\begin{equation}\label{eq:q-2}
U_{qq}(s_*,q)=\frac{S_*}{4\cos^2x\sqrt D}
                  +\frac{S_*\sin^2x}{D^{3/2}}>0.
\end{equation}
The endpoint at the lower branch is regular and has value
$1/\sin b$. Compactness supplies a positive lower bound on the whole
closed real $q$-interval. A sufficiently thin fixed convex rectangle
about that interval retains $\Rea U_{qq}\ge c$, also for $s$ in a
small complex neighborhood of $s_*$. At $q=0$, \eqref{eq:q-1}, or the dispersion
identity, gives $U=w_B$ and $U_q=0$ identically in $s$.

For the remaining compact left arc the limiting phase is
$\varphi=-\ii\gamma$, $\gamma>0$. There $q$ has real part
$\theta+b$, so its image is injective, and
$\dd q/\dd w=1+g'(w)$ cannot vanish: $g'$ is purely imaginary.
The two endpoint slivers remain regular after a small extension. The
left image has negative real part, the right image positive real part,
and they meet only in the regular branch chart. A finite cover by
analytic inverse neighborhoods therefore gives one inverse along this
embedded compact arc. Shrinking the neighborhoods, then fixing the
small deformation and parameter neighborhood, gives the asserted
bounded jets. The right inverse is the specific branch of \eqref{eq:q-1}.
\end{proof}

Set
\[
 \phi(s,q)=q-U(s,q)+w_B(s),\qquad
 a(s,q)=\frac{U_q(s,q)}{\sinh(\ii\phi(s,q))}.
\]
The coefficient $a$ is analytic and bounded at $q=0$: both its numerator
and denominator have a simple zero there, and
$U_q=q/\sin b+O(q^2)$ at $s=s_*$. It and its fixed jets are bounded
on the complete one-sided inverse neighborhoods.

The complete kernel is analytic in these regular coordinates. Near
the double branch the even inverse in $\varphi$ cancels the factor
$\varphi_i+\varphi_j$ in its denominator. Elsewhere that denominator
is separated, since the only limiting pair-product zero on the
one-sided support is the double lower branch. In particular on
$\mathcal D_R^2$,
\begin{equation}\label{eq:q-3}
|K(q_i,q_j)|\le C|q_i-q_j|,
\end{equation}
and fixed jets are bounded. The diagonal zero holds in the single inverse coordinate $q$ on
this domain, including its angular endpoint.

\subsection{Square-root geometry on the physical contour}
\begin{lemma}[Branch normal form]\label{lem:constrained-branch-geometry}
Put $u=\theta+b$ and $a_B=\sin b>0$.  On the original unit
contour, and on a true lower branch displaced outward by a common
$t_0\ge0$, set $E_0=\varepsilon+t_0$.  Uniformly for the fixed
small branch arc and $0\le t_0\le\tau$, one has
\[
 \begin{aligned}
 D&=1-W=a_Bu-i\beta+R,\qquad cE_0\le\beta\le CE_0,\\
 R&\in\mathbb R,\quad |R|\le C(u^2+E_0^2),\qquad
 D_u=a_B+O(|u|+E_0).
 \end{aligned}
\]
Here $\beta$ may depend smoothly on $u$, with
$|\beta_u|\le CE_0$.  On the physical fourth-quadrant branch,
\begin{equation}\label{eq:constrained-branch-bounds}
 \varphi=\sqrt{2D}\,h(D),\qquad h(0)=1,
 \qquad |\varphi'|\asymp(|u|+E_0)^{-1/2},
\end{equation}
where $h$ is analytic.  Moreover
\[
 -\Im\varphi\ge c\sqrt{|u|+E_0}\quad(u\le0),\qquad
 \Re\varphi'\ge\tfrac23|\varphi'|\quad(u\ge0),\qquad
 \int_{|u|\le\delta_B}|\varphi'(u)|\,du\le C.
\]
For coordinates with the same outward shift, a minimum nonnegative
coordinate $u_j$ satisfies $|\varphi'_i|\le C|\varphi'_j|$
for every other nonnegative coordinate.

\end{lemma}
\begin{proof}
Use $y=e^{t_0+i(-b+u)}$ and the exact real and imaginary parts of
$W=S-\cosh(t_0+i(-b+u))$.  The imaginary part is
\[
 \beta=\Im S+\sinh(t_0)\sin(b-u),
\]
and $\Im S=2\sin(\arg s_*)\varepsilon+O(\varepsilon^2)$.
The real part has the asserted expansion because
$1-S_*+\cos b=0$, $\Re S-S_*=O(\varepsilon^2)$, and
$\cosh t_0-1=O(t_0^2)$.  The analytic factorization of
$\arccos(1-D)$ gives \eqref{eq:constrained-branch-bounds}.
The relative error with respect to $a_Bu-i\beta$ is uniformly
small after the branch width and $\tau$ are fixed small.  Its
square root lies in a small fixed enlargement of the fourth quadrant.
For $u\le0$
its imaginary part has size at least a fixed multiple of its modulus;
for $u\ge0$ the derivative lies in a cone whose angle is at most
$\pi/4+o(1)$ from the positive real axis.  This proves the stated
attenuation and $2/3$ cone bounds.  The integrable majorant and the
minimum-coordinate comparison follow immediately from the two-sided
derivative estimate.

\end{proof}

\subsection{Coordinate size and angular diameter}

\begin{lemma}\label{lem:coordinate-scale}
On the actual right contour put
$u_i=\theta_i+b$, $t=N^{-1}\sum_i u_i$ and
$d=\max_i\theta_i-\min_i\theta_i>0$.
For any $r$ with $r\le d\le2r$, set $T=|t|+r+\varepsilon$.
Then
\begin{equation}\label{eq:q-21}
 |1+\varphi_w(w_i)|\le C(|u_i|+\varepsilon)^{-1/2},
\end{equation}
\begin{equation}\label{eq:q-22}
 \int_{t-2r}^{t+2r}(|u|+\varepsilon)^{-1/2}\dd u
       \le\frac{Cr}{\sqrt T},
\end{equation}
and
\begin{equation}\label{eq:q-23}
 |q_i|\le C\sqrt T,\qquad
 |q_i-q_l|+|K_{il}|\le Cr/\sqrt T.
\end{equation}
These are pointwise estimates. No extrema, diameter, or mean scale
is differentiated, and no partition of the integral is made here.
\end{lemma}
\begin{proof}
Near the lower branch, $p=0$, $m=1$, so the outward displacement
is the same $\tau P/N$ at every such coordinate. The branch normal
form gives \eqref{eq:q-21}; away from that branch compactness gives
the same bound, with a larger constant. It also holds along a
one-variable angular curve with the actual occupancies held fixed.
The curve derivative has the additional bounded factor
$1-\ii\tau[(P/N)m'-(Q/N)p']=1+O(\tau)$.

If $|t|\ge4r$, the denominator in \eqref{eq:q-22} is comparable
to $\sqrt{|t|+r+\varepsilon}$. Otherwise enlarge the interval to
$[-6r,6r]$ and use
\[
 \sqrt{6r+\varepsilon}-\sqrt\varepsilon
 =\frac{6r}{\sqrt{6r+\varepsilon}+\sqrt\varepsilon}.
\]
This proves the interval estimate even when the interval crosses
the branch. Joining two actual angles along their fixed-occupancy
curve gives exactly the actual $q$ endpoints. Integrating its
derivative, and then using \eqref{eq:q-3}, proves the second
inequality in \eqref{eq:q-23}.

To prove the absolute coordinate bound, first suppose $P=0$.
Nonnegativity of $p$ makes every $p(\theta_i)=0$, hence every
$h_i=0$. The branch normal form and compactness give
$|\varphi_i|\le C\sqrt{|u_i|+\varepsilon}$.
Since $q_i=\varphi_i+u_i-b-w_B(s)$ and
$w_B(s)+b=O(\varepsilon)$, this implies
$|q_i|\le C\sqrt{|u_i|+\varepsilon}\le C\sqrt T$.
If $P>0$, some actual right angle lies in the positive support
of $p$. Therefore $\max u_i\ge b$ and $|t|+2r\ge b$, so
$T\ge b/2$. Boundedness of the fixed regular domain gives the
same conclusion. All estimates allow arbitrary actual occupancy
means $P/N,Q/N$ in $[0,1]$ on the right rectangle.
\end{proof}

\section{Tangential differentiation of the complete one-sided integral}\label{sec:constrained-transport}\label{sec:global-transport}\subsection{The scalar density and its full Jacobian}

Let $M=D_\theta w=I-\ii Dh$. Its inverse is uniformly bounded for
small $\tau$, $\det M$ has bound $C^N$, and fixed jets add polynomial
factors in $N$. Also $\boldsymbol1^TM=\boldsymbol1^T$, hence
$\boldsymbol1^TM^{-1}=\boldsymbol1^T$. For
$\chi_s(\theta)=(q_i)$,
\begin{equation}\label{eq:q-4}
D_\theta\chi_s=\operatorname{diag}(1/U_{q,i})M,
 \quad J=\det D_\theta\chi_s
       =\det M\prod_i(1/U_{q,i}).
\end{equation}
This identity is used at $\varepsilon>0$, where all factors are
nonzero. The apparent branch poles will be retained as integrable
Jacobian factors, never differentiated as unregularized amplitudes.

The complete deformed scalar density is $J f(s,q)$, where
\begin{equation}\label{eq:q-5}
f(s,q)=\frac1{N!(2\pi)^{N-1}}
 \frac{1+Z}{1-Z}\det K\prod_i a(s,q_i).
\end{equation}
Formula \eqref{eq:q-5} includes the unchanged factor $2\pi$ converting the
ambient normalized density to constrained Haar normalization. The
entire occupancy Jacobian has cancelled from the regular amplitude.

For a residual field $B$ in $q$ coordinates define
\begin{equation}\label{eq:q-6}
V=-M^{-1}(\dot U+U_q B),
\end{equation}
where a dot denotes the fixed-$q$ derivative $\partial_s$ and products
inside the vector are componentwise. Differentiating
$U(s,\chi_s(\theta))=w(\theta)$ shows that \eqref{eq:q-6} is exactly
$(D_\theta\chi_s)^{-1}(\partial_s\chi_s-B)$.
If
\begin{equation}\label{eq:q-7}
\sum_i U_{q,i}B_i=-\sum_i\dot U_i,
\end{equation}
then $\sum_iV_i=0$ on an ambient neighborhood. If in addition
\begin{equation}\label{eq:q-8}
\sum_iB_i=-Nw_B'(s),
\end{equation}
then $(\partial_s-V\cdot\nabla_\theta)\sum_i\varphi_i=0$.
Indeed $\sum\varphi_i=\sum q_i-\sum w_i+Nw_B(s)$, and
$\sum w_i=\sum\theta_i$ is frozen by tangency.

For every fixed smooth real angular weight $\omega$ the exact identity
is
\begin{equation}\label{eq:q-9}
\partial_s(\omega Jf)-\operatorname{div}_\theta(V\omega Jf)
 =\omega J[\partial_sf+\operatorname{div}_q(Bf)]
    -(V\cdot\nabla_\theta\omega)Jf.
\end{equation}
To verify the identity, put $A_\chi=D_\theta\chi_s$.
Commuting mixed derivatives of $\chi_s$ gives
\[
 (\operatorname{div}_qB)\circ\chi_s
 =\operatorname{tr}\!\left[A_\chi^{-1}
   \left(\partial_sA_\chi-\sum_iV_i\partial_iA_\chi\right)\right]
     -\operatorname{div}_\theta V.
\]
Jacobi's formula therefore gives
$J(\operatorname{div}_qB)\circ\chi_s
=\partial_sJ-\operatorname{div}_\theta(JV)$.
The product and chain rules prove \eqref{eq:q-9}. Because $\sum V_i=0$ identically, the intrinsic divergence on
$\sum\theta_i=c$, parametrized by the first $N-1$ angles, is
\[
 \sum_{i<N}(\partial_i-\partial_N)V_i=\sum_i\partial_iV_i.
\]
The same formula holds for $V$ multiplied by any scalar density.
Thus ambient divergence restricts to constrained divergence without
an additional measure factor. Integration by parts applies to the
real and imaginary parts of these complex fields.
The identity is for scalar coefficients; it does not assert a new
intrinsic complex measure on $\mathbb H_N$. Smooth real bumps are
never analytically continued. Formal $q$ differentiation is the
commuting complex frame $(D_\theta\chi_s)^{-1}\nabla_\theta$; its
coefficients and their derivatives are controlled below.

More explicitly, if $A_\chi=D_\theta\chi_s$, set
\[
 E_a=\sum_i(A_\chi^{-1})_{ia}\partial_{\theta_i},\qquad
 E_s=\partial_s-
       \sum_i(A_\chi^{-1}\partial_s\chi_s)_i\partial_{\theta_i}.
\]
Then $E_a\chi_b=\delta_{ab}$ and $E_s\chi_b=0$. Every commutator
$[E_a,E_b]$ or $[E_s,E_a]$ is a linear combination of the real
angular derivatives and annihilates all $\chi_b$. Invertibility of
$A_\chi$ therefore makes every commutator zero. This is a pointwise
differential-operator calculation on the original real domain, and
requires no complex extension of its smooth occupancy functions.

\subsection{A bounded transport when a left root is present}\label{sec:left-gap}

Choose a single smooth right cutoff $r(\theta)$, equal to one for
$\theta+b\ge-\delta_B/2$ and zero for
$\theta+b\le-\delta_B$, with its harmless periodic return inside the
excluded upper plateau. Decompose the complete one-sided integral by
\[
 1=\prod_i r(\theta_i)+\left(1-\prod_i r(\theta_i)\right).
\]
Telescope the second term. Each of its at most $N$ terms has one fixed
named anchor $a$ in a compact left region, including on every cutoff
jet. That anchor has $|z_a|\le q_0<1$, and $|U_{q,a}|\ge c>0$.

On such a term take
\begin{equation}\label{eq:q-10}
B_a=-\frac{\sum_i\dot U_i}{U_{q,a}},\qquad B_i=0\ (i\ne a).
\end{equation}
Equation \eqref{eq:q-7} holds; equation \eqref{eq:q-8} is unnecessary because
$|1-Z|\ge1-q_0$. All fixed regular-coordinate and parameter jets of
$B$, $V$ and the regular amplitude have polynomial $N$ bounds, apart
from fixed one-body products costing $C_j^N$. In particular the
derivatives of $Z=\exp(-\ii\sum\phi_i)$ are bounded by a polynomial
in $N$, since $|Z|\le1$ at the actual contour and fixed jets of each
$\phi_i$ are bounded.

Hadamard applied to the complementary minors after fixed derivatives
therefore bounds the transported density by
\begin{equation}\label{eq:q-11}
C_j^N N^{-N/2+C_j}\prod_i |1+\varphi_w(w_i)|.
\end{equation}
The named compact factor is bounded. Eliminate its angle using the
constraint; every remaining factor has the independent majorant
$C(1+\operatorname{dist}(\theta,-b)^{-1/2})$. Integrating the product
costs $C^N$. This proves a uniform summable bound for every fixed
derivative of the whole left-root part, including all its mixed
branch configurations. No nested slope condition, product-phase
transport, coarea or convolution is needed for this piece.

\subsection{One collective field on the entire right part}\label{sec:right-field}

For this part all inverse coordinates lie in $\mathcal D_R$. Put
\[
 c_0=-w_B'(s),\quad
 A=-\sum_i\dot U_i-c_0\sum_iU_{q,i},\quad
 R=\sum_{i<j}(U_i-U_j)^2,
\]
and, off full equality,
\begin{equation}\label{eq:q-12}
B_i=c_0+\frac A R\sum_{j\ne i}
              \frac{(U_i-U_j)^2}{U_{q,i}-U_{q,j}}.
\end{equation}
Each summand extends analytically through its pair diagonal. Indeed
\[
 E(s;q,r)=\frac{U_q(s,q)-U_q(s,r)}{q-r}
         =\int_0^1U_{qq}(s,r+t(q-r))\,\dd t
\]
has $\Rea E\ge c$; hence, writing
$H=(U(s,q)-U(s,r))/(q-r)$, the summand is
$(q-r)H^2/E$. Pairing its antisymmetric terms gives
\begin{equation}\label{eq:q-13}
\sum_iB_i=Nc_0=-Nw_B',\qquad
 \sum_iU_{q,i}B_i=c_0\sum_iU_{q,i}+A=-\sum_i\dot U_i.
\end{equation}
Thus one field freezes both required quantities on the entire right
part, including every mixture of branch and compact coordinates.

\begin{lemma}[Uniform fixed-coordinate jets]\label{lem:jets}
Let $d=\max\theta_i-\min\theta_i$. On the actual right contour and
for any fixed total jet order,
\begin{equation}\label{eq:q-14}
|\partial_s^a\partial_q^\beta B_i|
  +|\partial_s^a\partial_q^\beta V_i|
 \le C_j N^{C_j}d^{-1-|\beta|}.
\end{equation}
There is no additional inverse power of $\varepsilon$.
\end{lemma}
\begin{proof}
First, at actual contour points,
$U_i-U_j=w_i-w_j=(\theta_i-\theta_j)-\ii(h_i-h_j)$.
The same occupancies occur in both $h_i,h_j$, so
\begin{equation}\label{eq:q-15}
|h_i-h_j|\le C\tau|\theta_i-\theta_j|,
 \quad \Rea R\ge(1-C^2\tau^2)\sum_{i<j}(\theta_i-\theta_j)^2.
\end{equation}
Consequently $|R|\ge c d^2$ and
$\sum_{i<j}|U_i-U_j|^2\le C|R|$.

Only a global diameter zero is needed for parameter derivatives.
The identities $U(s,0)=w_B(s)$ and $U_q(s,0)=0$ give, on the
fixed convex domain,
\begin{equation}\label{eq:q-16}
 |\partial_s^a U_q(s,q)|\le C_a|q|.
\end{equation}
Apply Lemma~\ref{lem:coordinate-scale} with $r=d$ and write
$T_d=|t|+d+\varepsilon$. Every point of the straight segment
joining $q_i$ and $q_j$ has modulus at most $C\sqrt{T_d}$.
Integration of \eqref{eq:q-16} along this segment gives
\begin{equation}\label{eq:q-17}
 |\partial_s^a(U_i-U_j)|
 \le C_a\sqrt{T_d}\,|q_i-q_j|\le C_a d.
\end{equation}
Consequently $|\partial_s^aR|\le C_aN^2d^2$.
Its first fixed-$q$ derivatives are $O(C_jN^{C_j}d)$, since
each differentiated product retains a pair difference.
Higher fixed coordinate derivatives are bounded by $C_jN^{C_j}$.
These assertions include arbitrary fixed additional parameter
derivatives, by \eqref{eq:q-17} and bounded analytic jets.
Together with $|R|\ge cd^2$, the quotient rule yields
\[
 |\partial_s^a\partial_q^\beta R^{-1}|
       \le C_jN^{C_j}d^{-2-|\beta|}.
\]
For $d\le1$, this follows by assigning a factor
$d^{2-|\beta|}$ to each jet of $R$; fixed parameter derivatives
cost no power. For $d\ge1$, the denominator is uniformly
separated and the fixed maximum diameter is absorbed in the constant.

Every fixed parameter jet of $H$ has size $O(\sqrt{T_d})$ by
\eqref{eq:q-16}, whereas all fixed jets of $E^{-1}$ are bounded.
Thus the pair numerator $F(q_i,q_j)=(q_i-q_j)H^2/E$ satisfies
\[
 |\partial_s^aF(q_i,q_j)|
 \le C_a|q_i-q_j|T_d\le C_a d\sqrt{T_d}\le C_a d.
\]
Every positive coordinate-order jet is bounded, including
additional fixed parameter derivatives, because $F$ is analytic
on the fixed product domain. It follows that
\[
 |\partial_s^a\partial_q^\beta\sum_{l\ne i}F(q_i,q_l)|
       \le C_jN^{C_j}d^{1-|\beta|}.
\]
All jets of $A$ are polynomially bounded in $N$.
Formula \eqref{eq:q-12} now proves the stated bound for $B$.

For the angular representative, the exact frame coefficients are
\[
 E_a\theta_i=(M^{-1})_{ia}U_{q,a},\qquad
 E_s\theta_i=(M^{-1}\dot U)_i.
\]
Their fixed jets are bounded by induction: differentiate
$M^{-1}$ by $D(M^{-1})=-M^{-1}(DM)M^{-1}$, and use the bounded
real angular derivatives of $M$ and the bounded analytic jets
of $U_q,\dot U$. Each fixed number of index sums costs only a
polynomial in $N$. The same induction applies to every original
smooth real cutoff and its fixed angular derivatives.
No inverse factor $U_q^{-1}$ is introduced. Applying the product
rule to \eqref{eq:q-6} therefore proves the same bound for $V$,
and justifies iterating \eqref{eq:q-9} with those real cutoffs.
\end{proof}

At fixed $\varepsilon,N$, the complete product identity gives
collision degree $N(N-1)$ in the original angular differences.
The field has one inverse diameter and each transport loses at most
two degrees. Thus the full-equality puncture flux vanishes through
order $j\le k$ whenever
\begin{equation}\label{eq:q-18}
N^2-2j-3>0,\qquad N\ge N_*+2,\quad k=N_*^2/2-1.
\end{equation}
No partial-diagonal flux occurs, since all pair quotients in \eqref{eq:q-12}
are regular. Perform the transports on the whole smooth right term
before selecting extrema or diameter shells.

For clarity, a derivative of a determinant marks one matrix entry;
its complementary minor retains Hadamard's bound. A fixed number
of derivatives has only polynomially many index choices. Each
field coefficient costs one inverse diameter, its coordinate
derivative costs one more, and its parameter derivative costs none.
The product rule and divergence therefore give, by induction,
the following filtration:
a descendant after $j$ transports contains $\ell$ marked entries and
a scalar loss $d^{-h}$ with
\begin{equation}\label{eq:q-19}
h+\ell\le2j.
\end{equation}
All regular entry, one-body and scalar jets are bounded with
polynomial $N$ costs. If the unchanged entries are at most $a$,
Hadamard and $N!\ge(N/e)^N$ leave
\begin{equation}\label{eq:q-20}
C_j^N N^{-N/2+C_j}d^{-h}a^{N-\ell}
 \frac{\prod_i|1+\varphi_w(w_i)|}{|1-Z|}.
\end{equation}

\section{One coarea estimate for the complete right integral}\label{sec:global-integration}
\subsection{A projected periodic kernel}
\begin{lemma}[Oblique periodic kernel]\label{lem:oblique-kernel}
Fix $|\vartheta|<\pi/2$ and put $P_\vartheta=2\pi\cos\vartheta$.
If $G=x-ia$, $a\ge a_0$, and $0<a_0\le1$, then, with
$v=\Re(e^{-i\vartheta}G)$,
\[
 \frac1{|1-e^{-iG}|}
 \le \frac{C_\vartheta}
 {a_0+\operatorname{dist}(v,P_\vartheta\mathbb Z)}.
\]
Consequently, on any injective $v$-parametrized arc whose
$v$-image has length $L$, the integral of this majorant in $v$
is at most $C_\vartheta(1+L)(1+|\log a_0|)$.
\end{lemma}
\begin{proof}
For $0<a\le1$, use
\[
 |1-e^{-a-ix}|^2=(1-e^{-a})^2+2e^{-a}(1-\cos x)
 \ge c\{a^2+\operatorname{dist}(x,2\pi\mathbb Z)^2\}.
\]
It follows that
$|1-e^{-iG}|\ge c\{a+\operatorname{dist}(x,2\pi\mathbb Z)\}$.
For an integer $n$ minimizing the latter distance,
\[
 |v-2\pi n\cos\vartheta|
 \le\cos\vartheta\,|x-2\pi n|+|\sin\vartheta|a.
\]
Thus $a_0+\operatorname{dist}(v,P_\vartheta\mathbb Z)$ is
bounded by a constant times the preceding lower bound.
For $a\ge1$, $|1-e^{-iG}|\ge1-e^{-1}$ and the proposed
denominator is bounded above by $1+P_\vartheta/2$.
Integration over translates of one period proves the final claim.
\end{proof}


For the actual contours the attenuation floor is
$-\log|Z|\ge cN\varepsilon$. Indeed $\Im W_i\ge c\varepsilon$,
and $W_i=\cos\varphi_i$ remains in a bounded set. Writing
$\varphi_i=x_i-ia_i$ gives
$\Im W_i=\sin x_i\sinh a_i\le C a_i$, hence
$a_i\ge c\varepsilon$. Thus Lemma~\ref{lem:oblique-kernel}
costs at most $C(1+L)(1+|\log\varepsilon|)$ on an injective
projected phase interval of length $L$.

\subsection{Oblique extreme coarea with the full occupancy dependence}

Let $I=[-b-\delta_B,\delta_R]$, where $\delta_R=\delta<b/2$,
and choose the negative segment inside the lower branch plateau.
The upper inverse branch is outside $I$.  Use the original balanced
contour with reflected disjoint bumps and put
\[
 \sigma=P/N,\quad\upsilon=Q/N,\quad
 f(u;\sigma,\upsilon)
 =-b+u-i\tau\{\sigma m(-b+u)-\upsilon p(-b+u)\}.
\]
Here $u\in[-\delta_B,b+\delta_R]$ and the original one-sided
occupancies satisfy $P\le1/2$, $0\le Q\le N$.
The auxiliary analytic estimates in terms of $E=\varepsilon+\tau\sigma$
hold on the fixed box $0\le\sigma,\upsilon\le1$.
The final bound in Lemma~\ref{lem:global-coarea} concerns actual
occupancies computed from the angles, including on the enlarged
rectangle used below.
Let $\varphi_s(w)=\arccos(S(s)-\cos w)$ on the physical sheet,
$b_s(w)=\partial_w\varphi_s(w)$. Use the regular coordinate
defined above.
The regularized phase coordinate is
$q_s(w)=\varphi_s(w)+w-w_B(s)$.
Its analytic inverse and the scalar transport were proved in
Sections~\ref{sec:global-coordinate}--\ref{sec:global-transport}.

\begin{lemma}[Occupancy-dependent extreme derivative]\label{lem:global-coarea}
Write
$G(\theta)=\sum_i\varphi_s(w_i(\theta))$ and
$\vartheta=3\pi/8$.  On an actual extreme-index cell with
$0\le m=\min_i u_i<M=\max_i u_i$, all other $u_i$ fixed,
and $m+M$ fixed by the constraint, one has
\begin{equation}\label{eq:g-G1}
\Re\left(e^{-i\vartheta}\frac{dG}{dm}\right)>0,
 \qquad
 \frac{|(b_s(w_m)+1)(b_s(w_M)+1)|}
 {\Re(e^{-i\vartheta}dG/dm)}
 \le C\frac{\sqrt{M+\varepsilon}}{M-m}.
\end{equation}
The constants are independent of $N$, $\varepsilon$, and the
actual occupancies.  A sufficiently small fixed $\tau$ is chosen
after the bump derivatives and fixed angular margins.
\end{lemma}
\begin{proof}
Temporarily hold the actual values of $\sigma,\upsilon$ fixed
and set $E=\varepsilon+\tau\sigma$.
On the plateau $f=-b+u-i\tau\sigma$; the square-root normal form gives
\begin{equation}\label{eq:g-G2}
|b_s(f(u))|\le C(u+E)^{-1/2},\qquad
 |\varphi_s''(f(u))|\asymp(u+E)^{-3/2},\qquad u\ge0.
\end{equation}
These estimates extend to the remainder of the fixed right
interval by compactness: it stays a fixed distance from both
dispersion branches, while $u$ is bounded above and below there.
The exact identity
\[
 -\varphi_s''(w)
 =\frac{S(s)(1-S(s)\cos w+\cos^2w)}{\sin^3\varphi_s(w)}
\]
also gives
\begin{equation}\label{eq:g-G3}
\Re\{e^{-i\vartheta}[-\varphi_s''(f(u))]\}
 \ge c|\varphi_s''(f(u))|.
\end{equation}
Indeed on the true branch the normal form puts
$\sin\varphi$ in an arbitrarily small enlargement of
$[-\pi/4,0]$.  The real branch-center displacement is
$O(\varepsilon^2+(\tau\sigma)^2)$ and the imaginary margin is
comparable to $E$, so this statement includes $u=0$.
On the compact remainder $\sin\varphi$ is close to a positive
real number.  The numerator is uniformly close to a positive
real number since
$1-S_*c+c^2\ge1-S_*^2/4>0$ for real $c$.
Thus \eqref{eq:g-G3} holds after fixing small $\tau,\varepsilon_0$.

Put $a(u)=\partial_u f(u)$ and $B(u)=b_s(f(u))a(u)$ at these
fixed occupancies.  Since $a=1+O(\tau)$,
$a'=O(\tau)$, and $|b_s(f)|\le C|\varphi_s''(f)|$ by \eqref{eq:g-G2},
\[
 -B'=-\varphi_s''(f)a^2-b_s(f)a'
\]
retains \eqref{eq:g-G3}.  Therefore, with $d=M-m$,
\begin{equation}\label{eq:g-G4}
\begin{aligned}
 &\Re\{e^{-i\vartheta}[B(m)-B(M)]\}\ge cL,\\
 &L=\int_m^M(u+E)^{-3/2}du
 =\frac{2d}{\sqrt{m+E}\sqrt{M+E}
 (\sqrt{m+E}+\sqrt{M+E})}.
\end{aligned}
\end{equation}

We now restore the entire occupancy dependence rather than
moving the occupancies along that auxiliary comparison curve.
Set
\[
 B_m^{\rm occ}=\sum_v m(\theta_v)b_s(w_v),\qquad
 B_p^{\rm occ}=\sum_v p(\theta_v)b_s(w_v).
\]
Direct differentiation of the original balanced contour gives
the exact formula
\begin{equation}\label{eq:g-G5}
\partial_{\theta_i}G
 =b_s(w_i)a(u_i)
 -\frac{i\tau p'(\theta_i)}N B_m^{\rm occ}
 +\frac{i\tau m'(\theta_i)}N B_p^{\rm occ}.
\end{equation}
The constraint slice differentiates in the direction
$\partial_{\theta_{\min}}-\partial_{\theta_{\max}}$.
Every actual $u_v\ge m$, and the true branches have the same
outward displacement $\tau\sigma$.  Hence \eqref{eq:g-G2} implies
\[
 \frac{|B_m^{\rm occ}|+|B_p^{\rm occ}|}{N}
 \le C(m+E)^{-1/2}.
\]
Both bump-derivative differences in \eqref{eq:g-G5} are bounded by $Cd$.
The total occupancy correction to $B(m)-B(M)$ is thus at most
$C\tau d(m+E)^{-1/2}$.
Since $M+E$ lies in a fixed bounded interval, the exact
expression in \eqref{eq:g-G4} gives
$d(m+E)^{-1/2}\le CL$.  Choosing $\tau$ small absorbs the
correction into half the positive projection in \eqref{eq:g-G4}.
This comparison uses the actual extreme configuration only;
it does not presume a Hessian bound along a transposition path
on which unchanged indices might become more extreme.

The numerator in \eqref{eq:g-G1} is at most
$C[(m+E)(M+E)]^{-1/2}$, with the constant enlarged on the bounded interval.
Combining with \eqref{eq:g-G4} bounds the ratio by $C\sqrt{M+E}/d$.
If $P=0$, then $E=\varepsilon$.  If $P>0$, an actual angle
belongs to the positive support of $p$, so $M\ge b$.
Therefore $E\le C(M+\varepsilon)$ in either event, proving
\eqref{eq:g-G1}.
\end{proof}



\paragraph{The geometric integration domain.}\label{par:right-rectangle}
After all differentiation, bound the smooth weight and its
transported coefficients by their proved envelopes and enlarge
the nonnegative integral to the full right angular rectangle
$I^N\cap\mathbb H_N$. All actual occupancies there still satisfy
$0\le\sigma,\upsilon\le1$, so the preceding lemma applies.
The balanced displacement has the same sign on each semicircle,
and hence $\Im(S-\cos w_i)\ge\Im S$: the physical sheet and
attenuation floor are preserved on this enlarged domain.
The regular inverse also remains valid there. On the compact
right arc its $q$-image changes by $O(\varepsilon+\tau)$;
near the branch, uniform square-root continuity bounds the
change by $O(\sqrt{\varepsilon+\tau})$. Choose $\delta_B$,
$\tau$ and $\varepsilon_0$ after the fixed margin of
$\mathcal D_R$ so that all these images lie inside it.

For fixed other coordinates and a fixed constrained sum,
the conditions specifying the extreme indices, a diameter
shell, and a nonnegative minimum are linear interval conditions
in the remaining variable $m$, with $M=A-m$.
Their intersection is therefore one interval. The negative
part is bounded directly below. No assertion about the number
of components of a smooth-cutoff fiber is required.

\subsection{A fixed mean scale for every diameter}


Write $u_i=\theta_i+b$, $t=N^{-1}\sum_i u_i$, and choose actual
minimum $m$ and maximum $M$ only in the final absolute integral.
The angular constraint fixes $t=b+2\pi l/N$; at most $CN$ such means
meet the right interval. The independent coordinate measure after
eliminating any one angle is unchanged.

On a shell $r\le d\le2r$ use the scale
$T=|t|+r+\varepsilon$ from Lemma~\ref{lem:coordinate-scale}.
Every coordinate belongs to $[t-2r,t+2r]$.

For a nonnegative minimum, Lemma~\ref{lem:global-coarea}
gives the monotone projected phase and
\begin{equation}\label{eq:q-24}
 \frac{|1+\varphi_w(w_m)|\,|1+\varphi_w(w_M)|}
 {|dv/dm|}\le C\frac{\sqrt{M+\varepsilon}}{M-m}.
\end{equation}
Its phase range is at most $CN$, since each phase is bounded.
Lemma~\ref{lem:oblique-kernel} therefore costs $CN(H+1)$.
Here and below $H=|\log\varepsilon|$.

For a nonnegative minimum, \eqref{eq:q-24} therefore bounds the integral in
the remaining extreme coordinate by
\begin{equation}\label{eq:q-25}
\int \frac{|1+\varphi_w(w_m)|\,|1+\varphi_w(w_M)|}{|1-Z|}\dd m
       \le CN\frac{\sqrt T}{r}(H+1).
\end{equation}
For a negative minimum, $m$ lies in the fixed true-branch segment,
where its attenuation gives
\[
 \frac{|1+\varphi_w(w_m)|\,|1+\varphi_w(w_M)|}{|1-Z|}
 \le \frac{C}{(|m|+\varepsilon)\sqrt{|M|+\varepsilon}}.
\]
Set $a=|m|+\varepsilon$, $b_1=|M|+\varepsilon$. Since
$a+b_1\ge M-m\ge r$, the elementary inequality
\begin{equation}\label{eq:q-26}
\frac1{a\sqrt{b_1}}
 \le \frac1{\sqrt r}\left(\frac1a+\frac1{b_1}\right)
\end{equation}
follows by multiplying through by $a\sqrt{b_1}$ and using
$a+b_1\ge r$. On the constrained slice $M=A-m$, so each reciprocal
in \eqref{eq:q-26} integrates to $C(H+1)$ over a bounded interval. This gives
\begin{equation}\label{eq:q-27}
\int \frac{|1+\varphi_w(w_m)|\,|1+\varphi_w(w_M)|}{|1-Z|}\dd m
 \le Cr^{-1/2}(H+1)\le C\frac{\sqrt T}{r}(H+1),
\end{equation}
which is the same pair envelope as \eqref{eq:q-25}. 

The two actual extreme Jacobian factors have been retained through
\eqref{eq:q-25}; only the $N-2$ remaining factors are now replaced by independent
one-body majorants. Their interval integrations cost
$(Cr/\sqrt T)^{N-2}$. Combine this with \eqref{eq:q-19}, \eqref{eq:q-20}, \eqref{eq:q-23}, and
\eqref{eq:q-25}--\eqref{eq:q-27}: the complete shell contribution is bounded by
\begin{equation}\label{eq:q-28}
C_j^NN^{-N/2+C_j}
 r^{2N-2j-3}T^{-N+3/2}(H+1)
 \le C_j^NN^{-N/2+C_j}r^{N-2j-3/2}(H+1).
\end{equation}
Here the fixed number of marked entries permits all bounded powers
of $T$ or $r$ to be absorbed into $C_j$. The last inequality uses
$T\ge r$. This estimate is valid for every $r>0$.

For all sufficiently large $N$, specifically $N>2k+4$, the exponent
in \eqref{eq:q-28} is positive for every $j\le k$. Summing dyadic shells and
the polynomial number of extrema and means gives a summable
$C_j^NN^{-N/2+C_j}(H+1)$ envelope.

For the finite orders $N_*+2\le N\le2k+4$, choose $\delta_B$ after
the full finite list $2\le n\le2k+4$, $n$ even, so that the distance
of every constrained mean from the branch exceeds $4\delta_B$.
The same choice will be used in the finite-order complementary lemma.
On the right interval $u\ge-\delta_B$,
a sufficiently small-diameter configuration must then have every
$u_i$ in a fixed compact right interval away from zero. There the
full fixed-$N$ Vandermonde degree gives the shell power
$N^2-2j-3>0$. The remaining shells have a fixed positive lower
diameter and \eqref{eq:q-28} suffices, even if its exponent is negative.
Thus the finite block has a uniform logarithmic bound as well.


\section{The complete logarithmic tail}
\label{sec:constrained-tail-sum}

\begin{theorem}\label{thm:constrained-tail}
At each selected point and for every $0\le j\le k$, uniformly
for $0<\varepsilon<\varepsilon_0$,
\begin{equation}\label{eq:logarithmic-absolute-tail}
 \sum_{\substack{N>N_*\\N\ \mathrm{even}}}
    |C_N^{(j)}((1+\varepsilon)s_*)|
       \le C_j(1+|\log\varepsilon|).
\end{equation}
In particular the entire higher-particle contribution is
$o(\varepsilon^{-1/2})$ at the derivative order of the first
singularity.
\end{theorem}
\begin{proof}
The selected part has the summable bound proved in
Lemma~\ref{lem:constrained-selected}. The left-root part has
the bound \eqref{eq:q-11}. For $N>N_c$, summing
\eqref{eq:q-28} over dyadic shells gives the same factorial
envelope times $H+1$. For the finitely many orders
$N_*+2\le N\le N_c$, the finite-mean argument following
\eqref{eq:q-28} gives a finite constant times $H+1$.
These estimates include every part of the exact identity
$C_N=F_N+2L_N$. Summing proves the theorem.
Since $\varepsilon^{1/2}(1+|\log\varepsilon|)\to0$, this is
the required noncancellation estimate.

The choices have the following order: the selected point,
$k$ and $N_c$; the fixed endpoint slivers and inverse-coordinate
neighborhood; the finite mean gap and a sufficiently small
branch width; the smooth bumps and right cutoff; a sufficiently
small fixed $\tau$; the protected parameter-disk constants;
and finally $\varepsilon_0$. Each estimate holds uniformly in
all tail orders after those choices. The diameter divisions
are made only after the complete weighted integrals have
been transported.
\end{proof}

\section{The first singularity on the constrained angular torus}
\label{sec:first-single}

Fix $\zeta=e^{i\alpha}$ from Lemma~\ref{lem:prime-root-family}, and put
\[
 N=N_*=2p,\qquad k=N^2/2-1,\qquad
 s_\varepsilon=(1+\varepsilon)\zeta,
 \qquad H=\log(1/\varepsilon).
\]
We use the complete coefficients $C_n$ in \eqref{eq:angular-bulk}.
The product-preserving deformation and its reflection decomposition
\eqref{eq:constrained-split} will also be used at the finitely many
orders $2\le n\le N$.  The formulas for the deformation and the
transport fields are independent of the order restriction used in
summing the infinite tail.

The geometric choices already made give $\delta<\alpha/2$ and
\eqref{eq:finite-mean-exclusion}. Hence $-\alpha$ lies on
$m=1,p=0$, $+\alpha$ lies on $p=1,m=0$, and no constrained
full-equality configuration of the present finite orders lies in
the chosen branch neighborhood.

\subsection{Finite-order complementary regions}

\begin{lemma}\label{lem:finite-complement}
Every positive even $n<N_*$ has bounded derivatives of every
fixed order along $s_\varepsilon$. At $N=N_*$, choose a fixed
smooth shape cutoff equal to one near
$\theta_1=\cdots=\theta_N=-\alpha$, supported where the
deformation is the identity and $w_L=1$. Let $I_N$ denote
that part of $L_N$. Then $C_N-2I_N$ has bounded derivatives
of every fixed order along this approach.
\end{lemma}
\begin{proof}
We first prove $O(1+H)$ for each fixed derivative. All
constants here may depend on the fixed order $n$ and the
number of derivatives.

For $F_n$, both occupancies are at least $1/4$. Each true
branch has displacement at least $\tau/(4n)$ and the product
has a positive damping gap at $s_*$. All one-body and pair
denominators are separated there: the only limiting possible
unit pair-product zero on these physical contours is at a
dispersion branch, and the displacements separate those
branches. Compactness supplies a fixed parameter disk on
which the localized integrand is bounded and holomorphic.
This argument includes $n=2$ and does not require the
fourfold convolution bound.

The left-root pieces of $L_n$ have bounded fixed derivatives
by the anchor construction of Section~\ref{sec:left-gap}.
Its argument is valid for each fixed $n\ge2$ and every
fixed number of derivatives: eliminate the named compact
anchor, retain its bounded Jacobian, and integrate the other
independent branch majorants. Thus only the complete right
part remains.

At full equality the constraint requires
$\theta_i=t_\theta=2\pi\ell/n$. By
\eqref{eq:finite-mean-exclusion}, none of these finitely many
points is at the branch. In the right interval their shifted
means satisfy $t_\theta+b>4\delta_B$, so the equality
neighborhoods are compact regular right charts. At equality
the balanced deformation is exactly zero, since
$P=np(t_\theta)$ and $Q=nm(t_\theta)$ give $h_i=0$.
Consequently its product denominator can vanish only if
\[
 n\arccos(S_*-\cos t_\theta)\in2\pi\mathbb Z,
 \qquad
 S_* = \cos(2\pi\ell/n)+\cos(2\pi m/n).
\]
Lemma~\ref{lem:prime-root-family} excludes this for $n<N_*$.
At $n=N_*$ it gives only $t_\theta=\pm\alpha$ modulo
$2\pi$; the positive point is outside the one-sided support.
Thus its only critical equality is $-\alpha$.

Localize each noncritical equality in a fixed small
neighborhood where the roots and pair factors are regular
and $1-Z$ is separated. Such localized integrals extend
holomorphically through $s_*$. At $N_*$ remove also the
specified chart at $-\alpha$. It lies entirely in $m=1,p=0$,
so its deformation is zero, $w_L=1$, and its right cutoff
is one. On the remaining right support, including every
cutoff derivative, the diameter satisfies $d\ge d_0>0$.

The single field of Section~\ref{sec:right-field} is now
smooth with bounded fixed regular jets, for any chosen
number of derivatives. There is no puncture flux. The
scalar transport gives the pointwise bound
\[
 C_{n,j}\frac{\prod_i|1+\varphi_w(w_i)|}{|1-Z|}.
\]
Retain the two actual extreme-coordinate factors for
Lemma~\ref{lem:global-coarea}; only after their absorption
are the other factors replaced by independent majorants.
For a negative minimum use \eqref{eq:q-26} instead.
There are finitely many means, extreme-index choices and
diameter shells with $d\ge d_0$. Equations
\eqref{eq:q-22}, \eqref{eq:q-25} and \eqref{eq:q-27}
therefore give $O(1+H)$ for every fixed derivative of this
entire remainder, including $n=2$ when there are no other
coordinates. The enlargement to a geometric interval
in Section~\ref{sec:global-integration} handles
disconnected smooth-cutoff fibers.

Each complete remainder is holomorphic on the exterior
sector: all contours and angular localizations are fixed
independently of $s$. For such a remainder $R$ and
$0<\varepsilon\le\varepsilon_1$,
\[
 R^{(j)}((1+\varepsilon)s_*)
 =R^{(j)}((1+\varepsilon_1)s_*)
   -s_*\int_\varepsilon^{\varepsilon_1}
              R^{(j+1)}((1+t)s_*)\,dt.
\]
The already proved bound at order $j+1$ is integrable at
zero, proving boundedness. Finally
$C_{N_*}-2I_{N_*}=F_{N_*}+2(L_{N_*}-I_{N_*})$ fixes the
reflection factor exactly.
\end{proof}

\subsection{The local shape integral}

Near the lower configuration write
\[
 \theta_i=-\alpha+t_i,\qquad
 t_N=-\sum_{i<N}t_i.
\]
The angular constraint is exactly $\sum_i t_i=0$, since
$N\alpha\in2\pi\mathbb Z$.  Its normalized measure is
$\prod_{i<N}dt_i/(2\pi)$.  All angles here are negative and lie on $m=1,p=0$,
so the deformation is identically zero and $w_L=1$.
The physical regular dispersion coordinate has the expansion
\[
 \varphi_i
  =\alpha-2i\varepsilon+t_i-\cot\alpha\,t_i^2
    +O(\varepsilon^2+\varepsilon|t_i|+|t_i|^3).
\]
Put $q(t,s)=1-Z(t,s)$.  Then
\begin{equation}\label{eq:single-q}
 \begin{split}
 q(t,s_\varepsilon)
  &=2N\varepsilon-i\cot\alpha\sum_i t_i^2
      +O(\varepsilon^2+\varepsilon|t|+|t|^3),\\
 \partial_s q(0,\zeta)&=2N\zeta^{-1}=:D_s.
 \end{split}
\end{equation}
The particle number is fixed in these expansions.

\begin{lemma}\label{lem:single-q-bound}
On a sufficiently small fixed real shape ball,
\[
 |q(t,s_\varepsilon)|\ge c(\varepsilon+|t|^2).
\]
\end{lemma}
\begin{proof}
Writing $\rho^2=\sum_i t_i^2$, the leading part in
\eqref{eq:single-q} satisfies
\[
 |2N\varepsilon-i\cot\alpha\,\rho^2|
       \ge c_0(\varepsilon+\rho^2).
\]
On $0<\varepsilon<\varepsilon_0$ and $\rho<\delta_0$, its
remainder is at most
$C(\varepsilon_0+\delta_0)(\varepsilon+\rho^2)$.
Choose the two fixed radii so that this is at most half the
preceding lower bound, and apply the triangle inequality.
\end{proof}

The local integral can be written exactly as
\begin{equation}\label{eq:angular-local-integral}
 I_N(s)=\int_{\mathbb R^{N-1}}
        \frac{\chi(t)\Delta(t)^2 A(t,s)}{q(t,s)}
                    \,dt_1\cdots dt_{N-1},
\end{equation}
where $\chi$ is the fixed shape cutoff and $A$ is analytic
near $(0,\zeta)$.  Its value is
\begin{equation}\label{eq:angular-local-prefactor}
 A(0,\zeta)=A_\alpha
 =\frac{2^{1-N(N-1)}}{N!(2\pi)^{N-1}\sin^{N^2}\alpha}>0.
\end{equation}
This is the coefficient computed in
Section~\ref{sec:elementary-pfaffian}, including
the determinant sign, all one-body factors, $1+Z$, $1/N!$,
and the normalized constrained measure.
All parameter derivatives of $A$ preserve the displayed
Vandermonde zeros, since the real angular coordinates are
independent of $s$.

\begin{lemma}\label{lem:single-period}
Put
\[
 a_0=\frac{N^2-1}{2},\qquad d_\alpha=\cot\alpha,\qquad
 \mathcal C_N=\frac1{\sqrt N}
   \int_{\substack{\sum_i\omega_i=0\\|\omega|=1}}
                   \Delta(\omega)^2\,d\sigma(\omega)>0.
\]
Then, with $\tau_N=-\sum_{i<N}\tau_i$,
\begin{equation}\label{eq:single-period}
 \begin{split}
 J_\alpha
 &=\int_{\mathbb R^{N-1}}
  \frac{\Delta(\tau)^2}
       {(2N-i d_\alpha\sum_i\tau_i^2)^{k+1}}
                      \,d\tau_1\cdots d\tau_{N-1}\\
 &=\frac{\mathcal C_N}{2}(2N)^{-1/2}
       (-i d_\alpha)^{-a_0}B(a_0,1/2)\ne0.
 \end{split}
\end{equation}
Here $\arg(-i d_\alpha)=-\pi/2$.
\end{lemma}
\begin{proof}
The coordinate measure on the zero-sum hyperplane is Euclidean
measure divided by $\sqrt N$.  Polar coordinates give radial
numerator $R^{N^2-2}$.  Substitution $u=R^2$ gives the beta
integral for a positive real quadratic coefficient; continue that
coefficient through the fourth quadrant to $-i d_\alpha$.
The uniform majorant is a constant times
$R^{N^2-2}(1+R^2)^{-N^2/2}$, whose tail is $R^{-2}$.
The angular constant is positive because a zero-sum vector with
distinct coordinates has an open spherical neighborhood where
$\Delta^2>0$.  Every remaining factor is nonzero.
\end{proof}

\begin{theorem}\label{thm:single-first}
At the selected prime root,
\[
 C_N^{(k)}(s_\varepsilon)
   =\Lambda_\alpha\varepsilon^{-1/2}
          +o(\varepsilon^{-1/2}),
 \qquad
 \Lambda_\alpha=2A_\alpha(-1)^kk!D_s^kJ_\alpha\ne0.
\]
For $j<k$, $C_N^{(j)}(s_\varepsilon)$ remains bounded.
Every even $n<N$ has bounded derivatives of every fixed order
along this approach.
\end{theorem}
\begin{proof}
By Lemma~\ref{lem:finite-complement}, $C_N-2I_N$ has bounded
fixed derivatives.  Differentiate \eqref{eq:angular-local-integral}
$k$ times.  Every term with denominator power at most $k$ is
uniformly integrable, since its radial exponent at
$\varepsilon=0$ is at least
\[
 N^2-2-2k=0.
\]
The only term of possible size $\varepsilon^{-1/2}$ is
\[
 \int \frac{\chi(t)\Delta(t)^2
    (-1)^kk!(\partial_s q(t,s_\varepsilon))^k A(t,s_\varepsilon)}
       {q(t,s_\varepsilon)^{k+1}}\,dt.
\]
Set $t=\sqrt\varepsilon\,\tau$.  Its total scale exponent is
\[
 \frac{N(N-1)}2+\frac{N-1}{2}-(k+1)=-\frac12.
\]
Lemma~\ref{lem:single-q-bound} gives the integrable majorant
$C\Delta(\tau)^2/(1+|\tau|^2)^{k+1}$.
Dominated convergence, \eqref{eq:single-q}, and
\eqref{eq:angular-local-prefactor} give
$A_\alpha(-1)^kk!D_s^kJ_\alpha$ for the lower chart.
The factor two in $\Lambda_\alpha$ counts the two reflected
centers; the separate factor two in $A_\alpha$ comes from $1+Z=2$.
For $j<k$ every denominator power is at most $k$, so the same
radial integrability argument and the complementary lemma prove
boundedness.  The lower-order conclusion is already contained
in that lemma.
\end{proof}

The radial phase is also explicit: $k=2p^2-1\equiv1\pmod4$ implies that
$\zeta^k\Lambda_\alpha$ is a positive real multiple of
$e^{-i\pi/4}$.
\section{Noncancellation and the natural boundary}

Let
\[
 U(s)=\sum_{\substack{N\ge2\\N\ \mathrm{even}}}C_N(s)
             =\X(s)/\Mtwo(s).
\]
The normally convergent exterior series may be differentiated term
by term. At a selected prime root, the finitely many lower orders
have bounded $k$th derivatives by Lemma~\ref{lem:finite-complement}.
Theorem~\ref{thm:single-first} gives the first singular term, and
Theorem~\ref{thm:constrained-tail} controls the complete higher tail.
Consequently
\begin{equation}\label{eq:U-leading}
 U^{(k)}(s_\varepsilon)
       =\Lambda_\alpha\varepsilon^{-1/2}
                           +o(\varepsilon^{-1/2}),
 \qquad \Lambda_\alpha\ne0.
\end{equation}
Thus $U$ cannot extend holomorphically through $\zeta$.
Since $\Mtwo$ is analytic and nonzero at this nonaxis point,
the exact identity $\X=\Mtwo U$ transfers noncontinuation to
the full connected bulk function.

For completeness, its physical leading coefficient follows directly.
The radial function $g(\varepsilon)=U((1+\varepsilon)\zeta)$
has $g^{(k)}(\varepsilon)=O(\varepsilon^{-1/2})$.
Integrability gives a finite limit of $g^{(k-1)}$ at zero;
successive integration gives finite limits of all lower derivatives.
Leibniz's rule therefore yields
\[
 \X^{(k)}(s_\varepsilon)
     =\Mtwo(\zeta)\Lambda_\alpha\varepsilon^{-1/2}
                             +o(\varepsilon^{-1/2}).
\]
The factor $\beta_{\rm phys}(s)=\operatorname{arsinh}(s)/(2J)$
is locally analytic and nonzero at every selected point on the
branch continued from the physical interval. Multiplication by it
gives the same obstruction for $\chi$.

The selected prime roots are dense in the open first quadrant.
Evenness and conjugation give a dense set of obstructions in all
four quadrants. Every neighborhood of a point on the unit circle,
including an axis point, contains such an obstruction. Hence no
point of the circle admits a holomorphic continuation of the
exterior germ. This proves Theorem~\ref{thm:main}.

\begin{thebibliography}{9}
\bibitem{Yang}
C. N. Yang, The spontaneous magnetization of a two-dimensional
Ising model, \emph{Physical Review} \textbf{85} (1952), 808--816.
\href{https://doi.org/10.1103/PhysRev.85.808}{doi:10.1103/PhysRev.85.808}.

\bibitem{PalmerTracy}
J. Palmer and C. Tracy, Two-dimensional Ising correlations:
Convergence of the scaling limit, \emph{Advances in Applied
Mathematics} \textbf{2} (1981), 329--388.
\href{https://doi.org/10.1016/0196-8858(81)90010-5}{doi:10.1016/0196-8858(81)90010-5}.

\bibitem{TracyWidom}
C. A. Tracy and H. Widom, On the singularities in the susceptibility
expansion for the two-dimensional Ising model,
\emph{Journal of Statistical Physics} \textbf{156} (2014), 1125--1135.
\href{https://doi.org/10.1007/s10955-014-1061-4}{doi:10.1007/s10955-014-1061-4}.

\end{thebibliography}
\end{document}
